\subsection{Bases}

\begin{exercise}[2b1]
    Find all vector spaces that have exactly one basis.

    \remark{The proof of this exercise fails if $\fbb$ contains only two elements.
	In this case, an additional vector space with exactly one basis is the vector space of dimension one over $\fbb$, since the only nonzero scalar is $1$, and any nonzero vector $v$ in the vector space is a basis.}
\end{exercise}

\begin{sol}
    Since an infinite-dimensional vector space has infinitely many bases, we only need to consider finite-dimensional vector spaces.
	Let $V$ be a finite-dimensional vector space, and let $v_1, \ldots, v_n$ be a basis of $V$.
	Then any other basis of $V$ must have the same length as $v_1, \ldots, v_n$, which is $n$.
	If $n \geq 1$, consider $\alpha \in \fbb$ such that $\alpha \neq 0$ and $\alpha \neq 1$.
	Then the list $\alpha v_1, v_2, \ldots, v_n$ still spans $V$, since we may recover $v_1 = \alpha\inv(\alpha v_1)$.
	Moreover, the list $\alpha v_1, v_2, \ldots, v_n$ is linearly independent, since if
    \[
    	\beta_1(\alpha v_1) + \beta_2 v_2 + \cdots + \beta_n v_n = 0,
    \]
    then we may write
    \[
    	(\alpha \beta_1) v_1 + \beta_2 v_2 + \cdots + \beta_n v_n = 0.
    \]
    Since $v_1, \ldots, v_n$ is linearly independent, we must have $\alpha \beta_1 = 0$, which implies that $\beta_1 = 0$.
	Then we also have $\beta_2 = \cdots = \beta_n = 0$.
	Thus, the list $\alpha v_1, v_2, \ldots, v_n$ is linearly independent.
	Therefore, if $n \geq 1$, then $V$ has more than one basis.
	If $n = 0$, then $V = \set{0}$.
	The empty list spans $V$, and the empty list must be linearly independent, since there are no nonzero scalars to consider.
	Thus, the empty list is a basis of $V$.
	Moreover, this is the only basis of $V$, since $0$ is linearly dependent.
	Therefore, the only vector space that has exactly one basis is the zero vector space.
\end{sol}

\begin{exercise}[2b2]
    Verify all assertions in Example 2.27.
    \begin{subexercises}
        \item The list $(1, 0, \ldots, 0), (0, 1, 0, \ldots, 0), \ldots, (0, \ldots, 0, 1)$ is a basis of $\fbb^n$, called the \emph{standard basis} of $\fbb^n$.
        \item The list $(1, 2), (3, 5)$ is a basis of $\fbb^2$.
	Note that this list has length two, which is the same as the length of the standard basis of $\fbb^2$.
	In the next section, we will see that this is not a coincidence.
        \item The list $(1, 2, -4), (7, -5, 6)$ is linearly independent in $\fbb^3$ but is not a basis of $\fbb^3$ because it does not span $\fbb^3$.
        \item The list $(1, 2), (3, 5), (4, 13)$ spans $\fbb^2$ but is not a basis of $\fbb^2$ because it is not linearly independent.
        \item The list $(1, 1, 0), (0, 0, 1)$ is a basis of $\set{(x, x, y) \in \fbb^3 \mid x, y \in \fbb}$.
        \item The list $(1, -1, 0), (1, 0, -1)$ is a basis of
        \[
        	\set{(x, y, z) \in \fbb^3 \mid x + y + z = 0}.
        \]
        \item The list $1, z, \ldots, z^m$ is a basis of $\pcal_m(\fbb)$, called the \emph{standard basis} of $\pcal_m(\fbb)$.
    \end{subexercises}
\end{exercise}

\begin{sole}
    \item Suppose $\alpha_1, \ldots, \alpha_n \in \fbb$ such that
    \[
    	\alpha_1(1, 0, \ldots, 0) + \alpha_2(0, 1, 0, \ldots, 0) + \cdots + \alpha_n(0, \ldots, 0, 1) = 0.
    \]
    Then each $\alpha_i = 0$, so the list is linearly independent.
	Moreover, any vector $(x_1, \ldots, x_n) \in \fbb^n$ can be written as
    \[
    	(x_1, \ldots, x_n) = x_1(1, 0, \ldots, 0) + x_2(0, 1, 0, \ldots, 0) + \cdots + x_n(0, \ldots, 0, 1),
    \]
    so the list spans $\fbb^n$.
	Thus, the list is a basis of $\fbb^n$.
    \item If $(3, 5) = a(1, 2)$ for some $a \in \fbb$, then we must have $3 = a$ and $5 = 2a$, which is impossible.
	Thus, the list is linearly independent.
	Moreover, for any element $(x, y) \in \fbb^2$, we can solve the system of equations
    \begin{align*}
        x & = 3\alpha + \beta, \\
        y & = 5\alpha + 2\beta
    \end{align*}
    for $\alpha, \beta \in \fbb$.
	Solving this system of equations, we find that
    \[
    	\alpha = 2x - y, \quad \beta = -5x + 3y.
    \]
    Since we can find $\alpha, \beta \in \fbb$ for any $(x, y) \in \fbb^2$, the list spans $\fbb^2$.
	Thus, the list is a basis of $\fbb^2$.
    \item Note that $\fbb^3$ is spanned by the linearly independent list $(1, 0, 0), (0, 1, 0), (0, 0, 1)$.
	Since the length of a spanning list must be at least the length of a linearly independent list, the list $(1, 2, -4), (7, -5, 6)$ cannot span $\fbb^3$.
    \item From part (b), we saw that $(1, 2), (3, 5)$ is a basis of $\fbb^2$.
	Since $(4, 13) = -5(3, 5) + 19(1, 2)$, the list $(1, 2), (3, 5), (4, 13)$ is linearly dependent.
	Moreover, since $(1, 2), (3, 5)$ is a basis of $\fbb^2$, the list $(1, 2), (3, 5), (4, 13)$ spans $\fbb^2$.
	Thus, the list is not a basis of $\fbb^2$.
    \item Let $U$ be the subspace in question.
	If $\alpha, \beta \in \fbb$ such that
    \[
    	\alpha(1, 1, 0) + \beta(0, 0, 1) = (0, 0, 0),
    \]
    then we must have $\alpha = \beta = 0$, so the list is linearly independent.
	Moreover, for any $(x, x, y) \in U$, we can write
    \[
    	(x, x, y) = x(1, 1, 0) + y(0, 0, 1),
    \]
    so the list spans $U$.
	Thus, the list is a basis of $U$.
    \item Let $U$ be the subspace in question.
	If $\alpha, \beta \in \fbb$ such that
    \[
    	\alpha(1, -1, 0) + \beta(1, 0, -1) = (0, 0, 0),
    \]
    then we must have $\alpha = \beta = 0$, so the list is linearly independent.
	Moreover, for any $(x, y, z) \in U$, note that $x + y + z = 0$, implying that $x = -y - z$.
	Then we can write
    \[
    	(x, y, z) = (-y - z, y, z) = -y(1, -1, 0) - z(1, 0, -1),
    \]
    so the list spans $U$.
	Thus, the list is a basis of $U$.
    \item If $\alpha_0, \ldots, \alpha_m \in \fbb$ such that
    \[
    	\alpha_0 + \alpha_1 z + \cdots + \alpha_m z^m = 0,
    \]
    then we must have $\alpha_0 = \cdots = \alpha_m = 0$, so the list is linearly independent.
	Moreover, any polynomial $p(z) \in \pcal_m(\fbb)$ can be written as
    \[
    	p(z) = \alpha_0 + \alpha_1 z + \cdots + \alpha_m z^m,
    \]
    so the list spans $\pcal_m(\fbb)$.
	Thus, the list is a basis of $\pcal_m(\fbb)$.
\end{sole}

\begin{exercise}[2b3]
    \begin{subexercises}
        \item Let $U$ be the subspace of $\rbb^5$ defined by
        \[
        	U = \set{(x_1, x_2, x_3, x_4, x_5) \in \rbb^5 \mid x_1 = 3x_2 \text{ and } x_3 = 7x_4}.
        \]
        Find a basis of $U$.
        \item Extend the basis in (a) to a basis of $\rbb^5$.
        \item Find a subspace $W$ of $\rbb^5$ such that $\rbb^5 = U \op W$.
    \end{subexercises}
\end{exercise}

\begin{sole}
    \item Note that we can rewrite the subspace $U$ as
    \[
    	U = \set{(3x_2, x_2, 7x_4, x_4, x_5) \in \rbb^5 \mid x_2, x_4, x_5 \in \rbb}.
    \]
    Then we can write any vector in $U$ as a linear combination of the following three vectors:
    \[
    	u_1 = (3, 1, 0, 0, 0), \quad u_2 = (0, 0, 7, 1, 0), \quad u_3 = (0, 0, 0, 0, 1).
    \]
    Moreover, these three vectors are linearly independent, since the only way to write the zero vector as a linear combination of $u_1, u_2, u_3$ is to take all coefficients to be zero.
	Thus, $\set{u_1, u_2, u_3}$ is a basis of $U$.
    \item We can extend the basis $\set{u_1, u_2, u_3}$ of $U$ to a basis of $\rbb^5$ by adding two more vectors that are linearly independent of $u_1, u_2, u_3$.
	For example, we can add the vectors
    \[
    	u_4 = (1, 0, 0, 0, 0), \quad u_5 = (0, 0, 1, 0, 0),
    \]
    which are linearly independent of $u_1, u_2, u_3$: since $u_1$ is the only vector in the list with a nonzero first component, $u_4$ cannot be a linear combination of $u_1, u_2, u_3$.
	Similarly, since $u_2$ is the only vector in the list with a nonzero third component, $u_5$ cannot be a linear combination of $u_1, u_2, u_3, u_4$.
	Thus, $\set{u_1, u_2, u_3, u_4, u_5}$ is a basis of $\rbb^5$.
    \item Let $W$ be the subspace of $\rbb^5$ spanned by the vectors $u_4$ and $u_5$.
	Then we have
    \[
    	\rbb^5 = U \op W,
    \]
    since $U \cap W = \set{0}$ and every vector in $\rbb^5$ can be written as a sum of a vector in $U$ and a vector in $W$.
\end{sole}

\begin{exercise}[2b4]
    \begin{subexercises}
        \item Let $U$ be the subspace of $\cbb^5$ defined by
        \[
        	U = \set{(z_1, z_2, z_3, z_4, z_5) \in \cbb^5 \mid 6z_1 = z_2 \text{ and } z_3 + 2z_4 + 3z_5 = 0}.
        \]
        Find a basis of $U$.
        \item Extend the basis in (a) to a basis of $\cbb^5$.
        \item Find a subspace $W$ of $\cbb^5$ such that $\cbb^5 = U \op W$.
    \end{subexercises}
\end{exercise}

\begin{sole}
    \item Note that we can rewrite the subspace $U$ as
    \[
    	U = \set{(z_1, 6z_1, -2z_4 - 3z_5, z_4, z_5) \in \cbb^5 \mid z_1, z_4, z_5 \in \cbb}.
    \]
    Then we can write any vector in $U$ as a linear combination of the following three vectors:
    \[
    	u_1 = (1, 6, 0, 0, 0), \quad u_2 = (0, 0, -2, 1, 0), \quad u_3 = (0, 0, -3, 0, 1).
    \]
    Moreover, these three vectors are linearly independent, since the only way to write the zero vector as a linear combination of $u_1, u_2, u_3$ is to take all coefficients to be zero.
	Thus, $\set{u_1, u_2, u_3}$ is a basis of $U$.
    \item We can extend the basis $\set{u_1, u_2, u_3}$ of $U$ to a basis of $\cbb^5$ by adding two more vectors that are linearly independent of $u_1, u_2, u_3$.
	For example, we can add the vectors
    \[
    	u_4 = (1, 0, 0, 0, 0), \quad u_5 = (0, 0, 1, 0, 0),
    \]
    which are linearly independent of $u_1, u_2, u_3$: since $u_1$ is the only vector in the list with a nonzero first component, $u_4$ cannot be a linear combination of $u_1, u_2, u_3$.
	Similarly, $u_2$ and $u_3$ have nonzero third components, but any linear combination of $u_2$ and $u_3$ will have nonzero fourth and fifth components, so $u_5$ cannot be a linear combination of $u_1, u_2, u_3, u_4$.
	Thus, $\set{u_1, u_2, u_3, u_4, u_5}$ is a basis of $\cbb^5$.
    \item Let $W$ be the subspace of $\cbb^5$ spanned by the vectors $u_4$ and $u_5$.
	Then we have
    \[
    	\cbb^5 = U \op W,
    \]
    since $U \cap W = \set{0}$ and every vector in $\cbb^5$ can be written as a sum of a vector in $U$ and a vector in $W$.
\end{sole}

\begin{exercise}[2b5]
    Suppose $V$ is finite-dimensional and $U, W$ are subspaces of $V$ such that $V = U + W$.
	Prove that there exists a basis of $V$ consisting of vectors in $U \cup W$.
\end{exercise}

\begin{sol}
    Let $u_1, \ldots, u_m$ be a basis of $U$ and $w_1, \ldots, w_n$ a basis of $W$.
	Then the list $u_1, \ldots, u_m, w_1, \ldots, w_n$ spans $V$, since any vector in $V$ can be written as a sum of a vector in $U$ and a vector in $W$.
	We can extract a basis of $V$ from this list by removing any linearly dependent vectors.
	Thus, there exists a basis of $V$ consisting of vectors in $U \cup W$.
\end{sol}

\begin{exercise}[2b6]
    Prove or give a counterexample: If $p_0, p_1, p_2, p_3$ is a list in $\pcal_3(\fbb)$ such that none of the polynomials $p_0, p_1, p_2, p_3$ has degree 2, then the list is not a basis of $\pcal_3(\fbb)$.
\end{exercise}

\begin{sol}
    This is not true.
	Let $p_0 = 1, p_1 = x, p_2 = x^3, p_3 = x^2 + x^3$.
	Then none of the polynomials $p_0, p_1, p_2, p_3$ has degree 2.
	Moreover, the list $p_0, p_1, p_2, p_3$ is linearly independent, since the only way to write the zero polynomial as a linear combination of $p_0, p_1, p_2, p_3$ is to take all coefficients to be zero.
	Finally, noting that $x^2 = p_3 - p_2$, then for any polynomial $p(x) = a_0 + a_1 x + a_2 x^2 + a_3 x^3 \in \pcal_3(\fbb)$, we can write
    \[
    	p(x) = a_0 p_0 + a_1 p_1 + a_2 (p_3 - p_2) + a_3 p_2 = a_0 p_0 + a_1 p_1 + (a_3 - a_2)p_2 + a_2 p_3,
    \]
    which shows that the list spans $\pcal_3(\fbb)$.
	Thus, the list $p_0, p_1, p_2, p_3$ is a basis of $\pcal_3(\fbb)$.
\end{sol}

\begin{exercise}[2b7]
    Suppose $v_1, v_2, v_3, v_4$ is a basis of $V$.
	Prove that
    \[
    	v_1 + v_2, v_2 + v_3, v_3 + v_4, v_4
    \]
    is also a basis of $V$.

    \remark{The technique of transforming the scalars $\beta_1, \beta_2, \beta_3, \beta_4$ into the scalars $\alpha_1, \alpha_2, \alpha_3, \alpha_4$ is good to remember.
	You will later learn that this is called a \emph{change of basis}.}
\end{exercise}

\begin{sol}
    Let $\alpha_1, \ldots, \alpha_4 \in \fbb$ such that
    \[
    	\alpha_1(v_1 + v_2) + \alpha_2(v_2 + v_3) + \alpha_3(v_3 + v_4) + \alpha_4 v_4 = 0.
    \]
    Then we may write
    \[
    	\alpha_1v_1 + (\alpha_1 + \alpha_2)v_2 + (\alpha_2 + \alpha_3)v_3 + (\alpha_3 + \alpha_4)v_4 = 0.
    \]
    Since $v_1, v_2, v_3, v_4$ is a basis of $V$, we must have
    \[
    	\alpha_1 = 0, \quad \alpha_1 + \alpha_2 = 0, \quad \alpha_2 + \alpha_3 = 0, \quad \alpha_3 + \alpha_4 = 0.
    \]
    From the first equation, we have $\alpha_1 = 0$.
	Substituting this into the second equation gives us $\alpha_2 = 0$.
	Substituting this into the third equation gives us $\alpha_3 = 0$.
	Finally, substituting this into the fourth equation gives us $\alpha_4 = 0$.
	Thus, the list is linearly independent.
	Moreover, since $v_1, v_2, v_3, v_4$ is a basis of $V$, then any vector $v \in V$ can be written as
    \[
    	v = \beta_1 v_1 + \beta_2 v_2 + \beta_3 v_3 + \beta_4 v_4
    \]
    for some $\beta_1, \ldots, \beta_4 \in \fbb$.
	Using the expressions from proving the linear independence of the list, we may show that there exist $\alpha_1, \ldots, \alpha_4 \in \fbb$ such that
    \[
        \begin{cases}
            \beta_1 = \alpha_1, \\
            \beta_2 = \alpha_1 + \alpha_2, \\
            \beta_3 = \alpha_2 + \alpha_3, \\
            \beta_4 = \alpha_3 + \alpha_4.
        \end{cases}
    \]
    Then $v \in V$ may be written as
    \[
    	v = \alpha_1(v_1 + v_2) + \alpha_2(v_2 + v_3) + \alpha_3(v_3 + v_4) + \alpha_4 v_4,
    \]
    where
    \[
    	\alpha_1 = \beta_1, \quad \alpha_2 = \beta_2 - \beta_1, \quad \alpha_3 = \beta_3 - \beta_2 + \beta_1, \quad \alpha_4 = \beta_4 - \beta_3 + \beta_2 - \beta_1.
    \]
    Thus, the list spans $V$.
	Therefore, the list is a basis of $V$.
\end{sol}

\begin{exercise}[2b8]
    Prove or give a counterexample: If $v_1, v_2, v_3, v_4$ is a basis of $V$ and $U$ is a subspace of $V$ such that $v_1, v_2 \in U$ and $v_3 \notin U$ and $v_4 \notin U$, then $v_1, v_2$ is a basis of $U$.
\end{exercise}

\begin{sol}
    This is not true.
	Let $V = \rbb^4$ and $v_1, v_2, v_3, v_4$ be the standard basis of $\rbb^4$.
	Let $U$ be the subspace of $\rbb^4$ defined by
    \[
    	U = \spn(v_1, v_2, v_3 + v_4) = \set{(x, y, z, z) \in \rbb^4 \mid x, y, z \in \rbb}.
    \]
    Then $v_1, v_2 \in U$ and $v_3, v_4 \notin U$.
	However, $v_1, v_2$ does not span $U$, since $(0, 0, 1, 1) \in U$ but cannot be written as a linear combination of $v_1$ and $v_2$.
	Thus, $v_1, v_2$ is not a basis of $U$.
\end{sol}

\begin{exercise}[2b9]
    Suppose $v_1, \ldots, v_m$ is a list of vectors in $V$.
	For $k \in \set{1, \ldots, m}$, let
    \[
    	w_k = v_1 + \cdots + v_k.
    \]
    Show that $v_1, \ldots, v_m$ is a basis of $V$ if and only if $w_1, \ldots, w_m$ is a basis of $V$.
\end{exercise}

\begin{sol}
    By \exref{2a14}, we know that $v_1, \ldots, v_m$ is linearly independent if and only if $w_1, \ldots, w_m$ is linearly independent.
	By \exref{2a3}, we know that $v_1, \ldots, v_m$ spans $V$ if and only if $w_1, \ldots, w_m$ spans $V$.
	Thus, $v_1, \ldots, v_m$ is a basis of $V$ if and only if $w_1, \ldots, w_m$ is a basis of $V$.
\end{sol}

\begin{exercise}[2b10]
    Suppose $U$ and $W$ are subspaces of $V$ such that $V = U \op W$.
	Suppose also that $u_1, \ldots, u_m$ is a basis of $U$ and $w_1, \ldots, w_n$ is a basis of $W$.
	Prove that
    \[
    	u_1, \ldots, u_m, w_1, \ldots, w_n
    \]
    is a basis of $V$.
\end{exercise}

\begin{sol}
    Let $u_1, \ldots, u_m$ be a basis of $U$ and $w_1, \ldots, w_n$ be a basis of $W$.
	We claim that the list
    \[
    	u_1, \ldots, u_m, w_1, \ldots, w_n
    \]
    is a basis of $V$.
	Let $\alpha_1, \ldots, \alpha_m, \beta_1, \ldots, \beta_n \in \fbb$ such that
    \[
    	\alpha_1 u_1 + \cdots + \alpha_m u_m + \beta_1 w_1 + \cdots + \beta_n w_n = 0.
    \]
    Then we may write
    \[
    	\alpha_1 u_1 + \cdots + \alpha_m u_m = -\beta_1 w_1 - \cdots - \beta_n w_n.
    \]
    Since the left-hand side is in $U$ and the right-hand side is in $W$, and $U \cap W = \set{0}$, we must have
    \[
    	\alpha_1 u_1 + \cdots + \alpha_m u_m = 0 \quad \text{and} \quad \beta_1 w_1 + \cdots + \beta_n w_n = 0.
    \]
    Since $u_1, \ldots, u_m$ is a basis of $U$, we have $\alpha_1 = \cdots = \alpha_m = 0$.
	Similarly, since $w_1, \ldots, w_n$ is a basis of $W$, we have $\beta_1 = \cdots = \beta_n = 0$.
	Thus, the list is linearly independent.
	Moreover, since $V = U \op W$, then $v = u + w$ for some $u \in U$ and $w \in W$.
	Then we may write
    \[
    	u = \gamma_1 u_1 + \cdots + \gamma_m u_m \quad \text{and} \quad w = \delta_1 w_1 + \cdots + \delta_n w_n
    \]
    for some $\gamma_1, \ldots, \gamma_m, \delta_1, \ldots, \delta_n \in \fbb$.
	Hence,
    \[
    	v = u + w = \gamma_1 u_1 + \cdots + \gamma_m u_m + \delta_1 w_1 + \cdots + \delta_n w_n,
    \]
    which shows that the list spans $V$.
	Therefore, the list is a basis of $V$.
\end{sol}

\begin{exercise}[2b11]
    Suppose $V$ is a real vector space.
	Show that if $v_1, \ldots, v_n$ is a basis of $V$ (as a real vector space), then $v_1, \ldots, v_n$ is also a basis of the complexification $V_\cbb$ (as a complex vector space).

    \note{Please see \exref{1b8} for the definition of the complexification $V_\cbb$}.
\end{exercise}

\begin{sol}
    Let $v_1, \ldots, v_n$ be a basis of $V$.
	We claim that $v_1, \ldots, v_n$ is also a basis of $V_\cbb$.
	Let $\alpha_1, \ldots, \alpha_n \in \cbb$ such that
    \[
    	\alpha_1 v_1 + \cdots + \alpha_n v_n = 0.
    \]
    Then we may write each $\alpha_k = a_k + b_k i$ for some $a_k, b_k \in \rbb$.
	Then we have
    \[
    	(a_1 + b_1 i)v_1 + \cdots + (a_n + b_n i)v_n = 0.
    \]
    This implies that
    \[
    	(a_1 v_1 + \cdots + a_n v_n) + i(b_1 v_1 + \cdots + b_n v_n) = 0.
    \]
    Since the left-hand side is a sum of two vectors in $V$, and $V$ is a real vector space, we must have
    \[
    	a_1 v_1 + \cdots + a_n v_n = 0 \quad \text{and} \quad b_1 v_1 + \cdots + b_n v_n = 0.
    \]
    Since $v_1, \ldots, v_n$ is a basis of $V$, we have $a_1 = \cdots = a_n = 0$ and $b_1 = \cdots = b_n = 0$.
	Thus, $\alpha_1 = \cdots = \alpha_n = 0$, which shows that the list is linearly independent.
	Moreover, let $v \in V_\cbb$.
	Then we may write $v = u + w i$ for some $u, w \in V$.
	Since $v_1, \ldots, v_n$ is a basis of $V$, we may write
    \[
    	u = \gamma_1 v_1 + \cdots + \gamma_n v_n \quad \text{and} \quad w = \delta_1 v_1 + \cdots + \delta_n v_n
    \]
    for some $\gamma_1, \ldots, \gamma_n, \delta_1, \ldots, \delta_n \in \rbb$.
	Hence,
    \[
    	v = u + w i = (\gamma_1 + \delta_1 i)v_1 + \cdots + (\gamma_n + \delta_n i)v_n,
    \]
    which shows that the list spans $V_\cbb$.
	Therefore, the list is a basis of $V_\cbb$.
\end{sol}